As the functor \(\Gamma\) is exact, the same holds for the
\(H^{n}(\mathbf{G},\,)\). It will now suffice for us to prove:

\begin{lemma}{5.2.2}
\label{I.5.2.2}
For every \(\mathbf{P}\in\Ob(\mathbf{O}\text{-}\mathrm{Mod.})\), one has:
\[
H^{n}(\mathbf{G},\Hom(\mathbf{G},\mathbf{P}))=0
\quad\text{and}\quad
\underline{H}^{n}(\mathbf{G},\uHom(\mathbf{G},\mathbf{P}))=0,
\quad\text{for }n>0.
\]
\end{lemma}

It suffices for us to prove that \(C^{*}(\mathbf{G},\Hom(\mathbf{G},\mathbf{P}))\) and
\(\underline{C}^{*}(\mathbf{G},\uHom(\mathbf{G},\mathbf{P}))\) are homotopically
trivial in degrees \(>0\). It even suffices to do so for the second, the
corresponding result for the first being deduced from it by base
change.\nde{The original has been expanded in what follows.}
Now, one defines for every \(n\geqslant 0\) a morphism
\[
\sigma\colon
\underline{C}^{n+1}(\mathbf{G},\uHom(\mathbf{G},\mathbf{P}))
\longrightarrow
\underline{C}^{n}(\mathbf{G},\uHom(\mathbf{G},\mathbf{P}))
\]
as follows. Let \(f\in\underline{C}^{n+1}(\mathbf{G},\uHom(\mathbf{G},\mathbf{P}))\);
for every \(S\in\Ob(\mathcal{C})\) and \(g_{1},\ldots,g_{n}\in\mathbf{G}(S)\),
\(\sigma(f)(g_{1},\ldots,g_{n})\) is the element of
\(\Hom_{S}(\mathbf{G}_{S},\mathbf{P}_{S})\) defined by: for every \(S'\to S\)
and \(x\in\mathbf{G}(S')\),
\[
\sigma(f)(g_{1},\ldots,g_{n})(x)
=f(x,g_{1},\ldots,g_{n})(e)\in\mathbf{P}(S')
\]
(where \(e\) denotes the unit element of \(\mathbf{G}(S')\)). Then \(\sigma\) is a
homotopy operator in degrees \(>0\). Indeed, for every
\(g_{1},\ldots,g_{n+1}\in\mathbf{G}(S)\) and \(x\in\mathbf{G}(S')\) one has, on
the one hand:
\begin{align*}
\partial\sigma(f)(g_{1},\ldots,g_{n+1})(x)
&=f(xg_{1},g_{2},\ldots,g_{n+1})(e)\\
&\quad+\sum_{i=1}^{n}(-1)^{i}f(x,g_{1},\ldots,g_{i}g_{i+1},\ldots,g_{n+1})(e)
\\
&\quad+(-1)^{n+1}f(x,g_{1},\ldots,g_{n})(e),
\end{align*}
and on the other hand:
% typo?: in \(\sigma(\partial f)\), the summand has no \((e)\), unlike the neighbouring terms; kept as printed.
\begin{align*}
\sigma(\partial f)(g_{1},\ldots,g_{n+1})(x)
&=(xf(g_{1},g_{2},\ldots,g_{n+1}))(e)
-f(xg_{1},g_{2},\ldots,g_{n+1})(e)\\
&\quad+\sum_{i=1}^{n}(-1)^{i+1}f(x,g_{1},\ldots,g_{i}g_{i+1},\ldots,g_{n+1})
\\
&\quad+(-1)^{n+2}f(x,g_{1},\ldots,g_{n})(e),
\end{align*}
whence
\[
\bigl(\partial\sigma(f)+\sigma(\partial f)\bigr)(g_{1},\ldots,g_{n+1})(x)
=f(g_{1},\ldots,g_{n+1})(x),
\]
\ie \(\partial\sigma+\sigma\partial\) is the identity map of
\(\underline{C}^{n+1}(\mathbf{G},\uHom(\mathbf{G},\mathbf{P}))\), for every
\(n\geqslant 0\).

\begin{remark}{5.2.3}
\label{I.5.2.3}
\nde{This remark has been added.}
The hypothesis ``\(\mathcal{C}\) small'' is used only to ensure the existence
of the derived functors \(R^{n}H^{0}\) and \(R^{n}\underline{H}^{0}\). In all
cases, what precedes shows that the functors \(H^{n}(\mathbf{G},\,)\)
\resp \(\underline{H}^{n}(\mathbf{G},\,)\) form a cohomological functor,
effaceable in degrees \(>0\), hence they are the (right) satellite functors
of the left exact functor \(H^{0}(\mathbf{G},\,)\)
\resp \(\underline{H}^{0}(\mathbf{G},\,)\), in the sense of \cite{Gr57}, 2.2;
if \((\mathbf{G}\text{-}\mathbf{O}\text{-}\mathrm{Mod.})\) has enough injective
objects (which is the case if \(\mathcal{C}\) is small), they coincide with
the derived functors (\loccit 2.3).
\end{remark}

\par\addvspace{0.55\baselineskip}%
\noindent\textbf{5.3.}\ Cohomology of the \(G\)-\(\OS\)-modules.\ ---\ \label{I.5.3}%
Let \(S\) be a scheme, \(G\) an \(S\)-group and \(\mathcal{F}\) a quasi-coherent
\(G\)-\(\OS\)-module. One defines the cohomology groups of \(G\) with values in
\(\mathcal{F}\) by
\[
H^{n}(G,\mathcal{F})=H^{n}(\mathbf{h}_{G},\mathbf{W}(\mathcal{F})).
\]
(for the notations, \cf 4.6).

Suppose \(G\) is \emph{affine} over \(S\). Then, in view of proposition 4.6.4,
this cohomology is computed in the following way: \(H^{n}(G,\mathcal{F})\) is
the \(n\)-th homology group of the complex \(C^{*}(G,\mathcal{F})\) whose
\(n\)-th term is:
\[
C^{n}(G,\mathcal{F})
=\Gamma\Bigl(S,\,
\mathcal{F}\otimes
\underbrace{\mathcal{A}(G)\otimes\cdots\otimes\mathcal{A}(G)}_{n\text{ times}}\Bigr).
\]
If \(f\) \resp \(a_{i}\) is a section of \(\mathcal{F}\) \resp of
\(\mathcal{A}(G)\) over an open of \(S\), one has
% typo?: no \(\otimes\) between \(a_{2}\) and \(\cdots\); kept as printed.
\begin{align*}
\partial(f\otimes a_{1}\otimes\cdots\otimes a_{n})
&=\mu_{\mathcal{F}}(f)\otimes a_{1}\otimes a_{2}\otimes\cdots\otimes a_{n}\\
&\quad+\sum_{i=1}^{n}(-1)^{i}f\otimes a_{1}\otimes\cdots\otimes\Delta a_{i}\otimes\cdots\otimes a_{n}\\
&\quad+(-1)^{n+1}f\otimes a_{1}\otimes a_{2}\cdots\otimes a_{n}\otimes 1,
\end{align*}
where \(\Delta\colon\mathcal{A}(G)\to\mathcal{A}(G)\otimes\mathcal{A}(G)\) and
\(\mu_{\mathcal{F}}\colon\mathcal{F}\to\mathcal{F}\otimes\mathcal{A}(G)\)
describe the coalgebra structure of \(\mathcal{A}(G)\) and the comodule
structure of \(\mathcal{F}\). Let us note in passing that the cohomology of
\(G\) with values in \(\mathcal{F}\) therefore depends only on the comodule
structure of \(\mathcal{F}\), and in particular only on the \(S\)-monoid
structure of \(G\).

One has in particular
\[
H^{0}(G,\mathcal{F})=\Gamma(S,\mathcal{F}^{G}),
\]
where \(\mathcal{F}^{G}\), the \emph{sheaf of invariants} of \(\mathcal{F}\),
is defined as the sheaf whose sections over the open \(U\) of \(S\) are the
sections of \(\mathcal{F}\) over \(U\) whose inverse image in every \(S'\)
over \(U\) is invariant under \(G(S')\) (\cf 4.7.1.2).

\begin{theorem}{5.3.1}
\label{I.5.3.1}
Let \(S\) be an affine scheme, \(G\) an \(S\)-group affine and flat over \(S\).
The functors \(H^{n}(G,\,)\) are the derived functors of \(H^{0}(G,\,)\) on the
category of quasi-coherent \(G\)-\(\OS\)-modules.
\end{theorem}

\begin{proof}
\nde{The original has been modified, in order to show that the category
\((\mathrm{G}\text{-}\OS\text{-}\mathrm{Mod.q.c.})\) is abelian and has enough
injective objects. One may compare with \cite{Ja03}, Part~I, 3.3--3.4, 3.9,
4.2 and 4.14--4.16 (where one will take care that ``\(k\)-group scheme'' means
``affine \(k\)-group scheme'', \cf\loccit, 2.1).}
As \(G\) is affine and flat over \(S\), then, by 4.7.2.1, the category
\((\mathrm{G}\text{-}\OS\text{-}\mathrm{Mod.q.c.})\) is equivalent to the
category \((\mathcal{A}(G)\text{-}\mathrm{Comod.q.c.})\) of quasi-coherent
\(\mathcal{A}(G)\)-comodules on \(\OS\), and is therefore abelian. On the
other hand, \(\mathcal{A}(G)\) being a flat \(\OS\)-module, each functor
\(\mathcal{F}\mapsto\mathcal{F}\otimes_{\OS}\mathcal{A}(G)^{\otimes n}\) is
exact; as moreover \(S\) is affine, one obtains that \(C^{*}(G,\,)\) is an
exact functor on \((\mathrm{G}\text{-}\OS\text{-}\mathrm{Mod.q.c.})\).

Let us denote by \(\Delta\) \resp \(\eta\) the comultiplication \resp the
augmentation of \(\mathcal{A}(G)\). For every quasi-coherent \(\OS\)-module
\(\mathcal{P}\), one denotes
\(\mathrm{Ind}(\mathcal{P})=\mathcal{P}\otimes_{\OS}\mathcal{A}(G)\) equipped
with the \(\mathcal{A}(G)\)-comodule structure defined by
\[
\mathrm{id}_{\mathcal{P}}\otimes\Delta\colon
\mathcal{P}\otimes_{\OS}\mathcal{A}(G)
\longrightarrow
\mathcal{P}\otimes_{\OS}\mathcal{A}(G)\otimes_{\OS}\mathcal{A}(G);
\]
this defines a functor
\(\mathrm{Ind}\colon(\OS\text{-}\mathrm{Mod.q.c.})\to(\mathrm{G}\text{-}\OS\text{-}\mathrm{Mod.q.c.})\).

It follows from 4.6.4.1, 5.2.0 and 5.2.0.1 that one has an isomorphism of
\(G\)-\(\OS\)-modules:
\[
(*)\qquad
\mathbf{W}(\mathrm{Ind}(\mathcal{P}))
\simeq
E(\mathbf{W}(\mathcal{P}))
=\Hom(G,\mathbf{W}(\mathcal{P})).
\]
Via this identification, the morphism
\(\varepsilon\colon E(\mathbf{W}(\mathcal{P}))\to\mathbf{W}(\mathcal{P})\)
corresponds to the morphism of \(\OS\)-modules
\(\mathrm{id}_{\mathcal{P}}\otimes\eta\colon\mathrm{Ind}(\mathcal{P})\to\mathcal{P}\).

% typo?: \(\mathbf{W}\colon(\OS\text{-}\mathrm{Mod.})\to(\OS\text{-}\mathrm{Mod.})\) kept as printed.
One has already used that the functor
\(\mathbf{W}\colon(\OS\text{-}\mathrm{Mod.})\to(\OS\text{-}\mathrm{Mod.})\) is
fully faithful; the same holds, by definition 4.7.1, of its restriction to
\((\mathrm{G}\text{-}\OS\text{-}\mathrm{Mod.})\), \ie if \(\mathcal{M}\),
\(\mathcal{M}'\) are \(G\)-\(\OS\)-modules, one has a functorial isomorphism
\[
\Hom_{\mathrm{G}\text{-}\OS\text{-}\mathrm{Mod.}}(\mathcal{M},\mathcal{M}')
\simeq
\Hom_{\mathrm{G}\text{-}\OS\text{-}\mathrm{Mod.}}(\mathbf{W}(\mathcal{M}),\mathbf{W}(\mathcal{M}')).
\]
Consequently, one deduces from lemma 5.2.0.2 the

\begin{corollary}{5.3.1.1}
\label{I.5.3.1.1}
(i) The functor \(\mathrm{Ind}\) is right adjoint to the forgetful functor
\[
(\mathrm{G}\text{-}\OS\text{-}\mathrm{Mod.q.c.})\to(\OS\text{-}\mathrm{Mod.q.c.}).
\]
More precisely, the map
\[
\mathrm{id}_{\mathcal{P}}\otimes\eta\colon\mathrm{Ind}(\mathcal{P})\to\mathcal{P}
\]
induces, for every object \(\mathcal{M}\) of
\((\mathrm{G}\text{-}\OS\text{-}\mathrm{Mod.q.c.})\), a bijection
\[
\Hom_{\mathrm{G}\text{-}\OS\text{-}\mathrm{Mod.}}(\mathcal{M},\mathrm{Ind}(\mathcal{P}))
\xrightarrow{\sim}
\Hom_{\OS}(\mathcal{M},\mathcal{P}).
\]
(ii) Consequently, if \(\mathcal{I}\) is an injective object of
\((\OS\text{-}\mathrm{Mod.q.c.})\) then \(\mathrm{Ind}(\mathcal{I})\) is an
injective object of \((\mathrm{G}\text{-}\OS\text{-}\mathrm{Mod.q.c.})\).
\end{corollary}

Let \(\mathcal{F}\) be a \(G\)-\(\OS\)-module and
\(\mu_{\mathcal{F}}\colon\mathcal{F}\to\mathrm{Ind}(\mathcal{F})\) the map
defining the \(\mathcal{A}(G)\)-comodule structure. It follows from 5.2.0.3
(or else from the axioms (CM~1) and (CM~2) of 4.7.2) that \(\mu_{\mathcal{F}}\)
is a morphism of \(G\)-\(\OS\)-modules, and that
\((\mathrm{id}_{\mathcal{F}}\otimes\eta)\circ\mu_{\mathcal{F}}=\mathrm{id}_{\mathcal{F}}\),
hence that \(\mathcal{F}\) is a direct factor of \(\mathrm{Ind}(\mathcal{F})\)
as an \(\OS\)-module; in particular, \(\mu_{\mathcal{F}}\) is a monomorphism.
As one has, by \((*)\) and 5.2.2,
\[
H^{n}(G,\mathbf{W}(\mathrm{Ind}(\mathcal{F})))
\simeq
H^{n}(G,\Hom_{S}(G,\mathbf{W}(\mathcal{F})))
=0
\qquad\text{for }n>0,
\]
one therefore obtains that \(H^{n}(G,\,)\) is effaceable for \(n>0\).

Finally, \(S\) being affine, \((\OS\text{-}\mathrm{Mod.q.c.})\) has enough
injective objects. Let therefore \(\mathcal{F}\hookrightarrow\mathcal{I}\) be a
monomorphism of \(\OS\)-modules, where \(\mathcal{I}\) is an injective object
of \((\OS\text{-}\mathrm{Mod.q.c.})\); then, \(\mathcal{A}(G)\) being flat over
\(\OS\), \(\mathrm{Ind}(\mathcal{F})\) is a sub-\(G\)-\(\OS\)-module of
\(\mathrm{Ind}(\mathcal{I})\), whence:

\begin{corollary}{5.3.1.2}
\label{I.5.3.1.2}
The abelian category \((\mathrm{G}\text{-}\OS\text{-}\mathrm{Mod.q.c.})\) has
enough injective objects.
\end{corollary}

Taking into account \cite{Gr57}, 2.2.1 and 2.3 (already used in the proof of
5.2.1), this completes the proof of theorem 5.3.1.
\end{proof}

\begin{remark}{5.3.1.3}
\label{I.5.3.1.3}
One can also prove 5.3.1.1 by the following computation. To every morphism of
\(G\)-\(\OS\)-modules
\(\phi\colon\mathcal{M}\to\mathcal{P}\otimes_{\OS}\mathcal{A}(G)\) one
associates the \(\OS\)-morphism
\((\mathrm{id}_{\mathcal{P}}\otimes\eta)\circ\phi\colon\mathcal{M}\to\mathcal{P}\).
Conversely, to every \(\OS\)-morphism \(f\colon\mathcal{M}\to\mathcal{P}\) one
associates the morphism of \(G\)-\(\OS\)-modules
\((f\otimes\mathrm{id}_{\mathcal{A}(G)})\circ\mu_{\mathcal{M}}\colon\mathcal{M}\to\mathcal{P}\otimes_{\OS}\mathcal{A}(G)\).
One sees at once that
% typo?: \((f\otimes\mathrm{id}_{\OS})\circ(\mathrm{id}_{\mathcal{P}}\otimes\eta)\) kept as printed.
\begin{equation*}
(\mathrm{id}_{\mathcal{P}}\otimes\eta)\circ(f\otimes\mathrm{id}_{\mathcal{A}(G)})\circ\mu_{\mathcal{M}}
=(f\otimes\mathrm{id}_{\OS})\circ(\mathrm{id}_{\mathcal{P}}\otimes\eta)\circ\mu_{\mathcal{M}}
=f.
\end{equation*}

On the other hand, for every \(\phi\) the diagram below is commutative:
\[
\xymatrix{
\mathcal{M} \ar[r]^{\phi} \ar[d]_{\mu_{\mathcal{M}}}
  & \mathcal{P}\otimes_{\OS}\mathcal{A}(G) \ar[d]^{\mathrm{id}_{\mathcal{P}}\otimes\Delta} \\
\mathcal{M}\otimes_{\OS}\mathcal{A}(G) \ar[r]_-{\phi\otimes\mathrm{id}_{\mathcal{A}(G)}}
  & \mathcal{P}\otimes_{\OS}\mathcal{A}(G)\otimes_{\OS}\mathcal{A}(G).
}
\]
It follows that
\begin{align*}
\Bigl(\bigl((\mathrm{id}_{\mathcal{P}}\otimes\eta)\circ\phi\bigr)\otimes\mathrm{id}_{\mathcal{A}(G)}\Bigr)\circ\mu_{\mathcal{M}}
&=(\mathrm{id}_{\mathcal{P}}\otimes\eta\otimes\mathrm{id}_{\mathcal{A}(G)})\circ(\phi\otimes\mathrm{id}_{\mathcal{A}(G)})\circ\mu_{\mathcal{M}}\\
&=(\mathrm{id}_{\mathcal{P}}\otimes\eta\otimes\mathrm{id}_{\mathcal{A}(G)})\circ(\mathrm{id}_{\mathcal{P}}\otimes\Delta)\circ\phi
=\phi.
\end{align*}
This proves 5.3.1.1\,(i) (and (ii) follows).
\end{remark}

Let \(\mathcal{F}\) be a \(G\)-\(\OS\)-module; one has seen above that axiom (CM~2)
of 4.7.2 shows that, considered as an \(\OS\)-module, \(\mathcal{F}\) is a
direct factor of \(E(\mathcal{F})\). This entails:

\begin{proposition}{5.3.2}
\label{I.5.3.2}
Let \(S\) be an affine scheme and \(G\) an affine and flat \(S\)-group; suppose
that every exact sequence \(0\to\mathcal{F}_{1}\to\mathcal{F}_{2}\to\mathcal{F}_{3}\to 0\)
of quasi-coherent \(G\)-\(\OS\)-modules, which splits as a sequence of
\(\OS\)-modules, likewise splits as a sequence of \(G\)-\(\OS\)-modules.

Then the functors \(H^{n}(G,\,)\) are zero for \(n>0\) (or, what amounts to the
same, the functor \(H^{0}(G,\,)\) is \emph{exact}).
\end{proposition}

Indeed, by the hypothesis, the sequence of \(G\)-\(\OS\)-modules
\[
0\longrightarrow\mathcal{F}\longrightarrow E(\mathcal{F})\longrightarrow E(\mathcal{F})/\mathcal{F}\longrightarrow 0
\]
is split; \(\mathcal{F}\) is therefore a direct factor, as a
\(G\)-\(\OS\)-module, in \(E(\mathcal{F})\), but the cohomology of the latter is
zero.

One deduces immediately from this and from proposition 4.7.4:

\begin{theorem}{5.3.3}
\label{I.5.3.3}
Let \(S\) be an affine scheme and \(G\) a diagonalizable \(S\)-group. For every
quasi-coherent \(G\)-\(\OS\)-module \(\mathcal{F}\), one has
\(H^{n}(G,\mathcal{F})=0\), for \(n>0\).
\end{theorem}

\begin{remark*}
Proposition 5.3.2 remains valid when \(G\) is not necessarily flat over \(S\);
the proof then appeals to relative cohomology.\nde{The editors have not sought
to develop this remark.}
\end{remark*}

\sgasection{6}{\texorpdfstring{$G$-equivariant objects and modules}{G-equivariant objects and modules}}
\label{I.sec.6}
\nde{This section has been added.}
Let \(\mathcal{C}\) be a category having a final object and in which fiber
products exist. Let \(\mathbf{G}\) be a \(\widehat{\mathcal{C}}\)-group,
\(\pi\colon\mathbf{M}\to\mathbf{X}\) a morphism of \(\widehat{\mathcal{C}}\),
and \(\lambda=\lambda_{\mathbf{X}}\colon\mathbf{G}\times\mathbf{X}\to\mathbf{X}\)
an action of \(\mathbf{G}\) on \(\mathbf{X}\). In what follows, one will denote
by \(\mathbf{Y}\times_{f}\mathbf{M}\) the fiber product of
\(\pi\colon\mathbf{M}\to\mathbf{X}\) and of an \(\mathbf{X}\)-functor
\(f\colon\mathbf{Y}\to\mathbf{X}\).

For every \(U\in\Ob(\mathcal{C})\) and \(x\in\mathbf{X}(U)\), one will denote
\(\mathbf{M}_{x}=U\times_{x}\mathbf{M}\), \ie for every
\(\varphi\colon U'\to U\), one has
\[
\mathbf{M}_{x}(U')
=\{m\in\mathbf{M}(U')\mid \pi(m)=x_{U'}=\varphi^{*}(x)\}.
\]
Finally, if \(g\in\mathbf{G}(U)\) one will also denote by \(gx\) the element
\(\lambda(g,x)\) of \(\mathbf{X}(U)\).

\begin{definition}{6.1}
\label{I.6.1}
a) One says that \(\mathbf{M}\) is a \emph{\(G\)-equivariant
\(\mathbf{X}\)-object} if one has given oneself an action
\(\Lambda\colon\mathbf{G}\times\mathbf{M}\to\mathbf{M}\) of \(\mathbf{G}\) on
\(\mathbf{M}\) lifting \(\lambda\), \ie such that the square below is
commutative:
\[
\xymatrix{
\mathbf{G}\times\mathbf{M} \ar[r]^{\Lambda} \ar[d]_{\mathrm{id}_{\mathbf{G}}\times\pi}
  & \mathbf{M} \ar[d]^{\pi} \\
\mathbf{G}\times\mathbf{X} \ar[r]^{\lambda}
  & \mathbf{X}.
}
\]
This amounts to giving, for every morphism
\((g,x)\colon U\to\mathbf{G}\times\mathbf{X}\), maps
\[
\Lambda^{U}_{x}(g)\colon
\mathbf{M}_{x}(U)\longrightarrow\mathbf{M}_{gx}(U),
\qquad m\longmapsto g\cdot m
\]
satisfying \(1\cdot m=m\) and \(g\cdot(h\cdot m)=(gh)\cdot m\) and functorial in
the \((\mathbf{G}\times\mathbf{X})\)-object \(U\).
This amounts again to giving morphisms of \(U\)-objects:
\[
\Lambda_{x}(g)\colon
\mathbf{M}_{x}\longrightarrow\mathbf{M}_{gx}
\]
satisfying \(\Lambda_{x}(1)=\mathrm{id}\) and
\(\Lambda_{hx}(g)\circ\Lambda_{x}(h)=\Lambda_{x}(gh)\).

b) Let now \(\mathbf{O}\) be a \(\widehat{\mathcal{C}}\)-ring and let
\(\mathbf{O}_{\mathbf{X}}=\mathbf{O}\times\mathbf{X}\). Under the conditions of
(a), one says that \(\mathbf{M}\) is a \emph{\(G\)-equivariant
\(\mathbf{O}_{\mathbf{X}}\)-module} if it is an
\(\mathbf{O}_{\mathbf{X}}\)-module (\cf definition 4.3.3.1, valid for every
ring-valued functor on a category \(\mathcal{C}\)) and if the action
\(\Lambda\) is compatible with the \(\mathbf{O}_{\mathbf{X}}\)-module structure
of \(\mathbf{M}\), \ie if for every
\((g,x)\in\mathbf{G}(U)\times\mathbf{X}(U)\)
(\(U\in\Ob(\mathcal{C})\)), the map
\(\Lambda_{x}(g)\colon\mathbf{M}_{x}\to\mathbf{M}_{gx}\), \(m\mapsto g\cdot m\)
is a morphism of \(\mathbf{O}_{U}\)-modules.
\end{definition}

\begin{remark}{6.2}
\label{I.6.2}
(i) In (a) above, the conditions \(\Lambda_{x}(1)=\mathrm{id}\) and
\(\Lambda_{hx}(g)\circ\Lambda_{x}(h)=\Lambda_{x}(gh)\) evidently entail that
each \(\Lambda_{x}(g)\) is an \emph{isomorphism}, with inverse
\(\Lambda_{gx}(g^{-1})\). Conversely, if one supposes that each
\(\Lambda_{x}(g)\) is an isomorphism, the condition
\(\Lambda_{hx}(g)\circ\Lambda_{x}(h)=\Lambda_{x}(gh)\) applied to \(h=1\) gives
\(\Lambda_{x}(1)=\mathrm{id}\).

(ii) Let \(\mathbf{M}\) be an \(\mathbf{O}_{\mathbf{X}}\)-module. First, one
sees that giving a morphism \(\Lambda\colon\mathbf{G}\times\mathbf{M}\to\mathbf{M}\)
making the diagram of 6.1 commutative and such that each morphism
\(\Lambda_{x}(g)\colon\mathbf{M}_{x}\to\mathbf{M}_{gx}\), \(m\mapsto g\cdot m\),
is an isomorphism of \(\mathbf{O}_{U}\)-modules, amounts to giving an
\emph{isomorphism} of \(\mathbf{O}_{\mathbf{G}\times\mathbf{X}}\)-modules:
\begin{align*}
\theta\colon
\mathbf{G}\times\mathbf{M}
&=(\mathbf{G}\times\mathbf{X})\times_{\mathrm{pr}_{\mathbf{X}}}\mathbf{M}
\isomto
(\mathbf{G}\times\mathbf{X})\times_{\lambda}\mathbf{M}\\
&\qquad(g,x,m)\longmapsto(g,x,g\cdot m).
\end{align*}
As one has supposed that each \(\Lambda_{x}(g)\) is an isomorphism, the equality
\(\Lambda_{x}(1)=\mathrm{id}\) will be a consequence of the equality
\(\Lambda_{hx}(g)\circ\Lambda_{x}(h)=\Lambda_{x}(gh)\) applied to \(h=1\). One
therefore obtains that \(\Lambda\) is an \emph{action} of \(\mathbf{G}\) on
\(\mathbf{M}\) (\ie \(g\cdot(h\cdot m)=(gh)\cdot m\)) if and only if the
diagram of \((\mathbf{G}\times\mathbf{G}\times\mathbf{X})\)-isomorphisms below
is commutative (where one denotes by \(\mu\) the multiplication of
\(\mathbf{G}\) and by \(f^{*}(\theta)\) the isomorphism deduced from \(\theta\) by
a base change \(f\colon\mathbf{G}\times\mathbf{G}\times\mathbf{X}\to\mathbf{G}\times\mathbf{X}\)):
\[
\xymatrix@C=0.6em@R=2.6em{
(\mathbf{G}\times\mathbf{G}\times\mathbf{X})
  \times_{\mathrm{pr}_{\mathbf{X}}\circ\mathrm{pr}_{2,3}}
  \mathbf{M}
\ar[rr]^-{\mathrm{pr}_{2,3}^{*}(\theta)}_-{\sim}
&&
(\mathbf{G}\times\mathbf{G}\times\mathbf{X})
  \times_{\lambda\circ\mathrm{pr}_{2,3}}
  \mathbf{M}
\\
(\mathbf{G}\times\mathbf{G}\times\mathbf{X})
  \times_{\mathrm{pr}_{\mathbf{X}}\circ(\mu\times\mathrm{id}_{\mathbf{X}})}
  \mathbf{M}
\ar[d]_-{(\mu\times\mathrm{id}_{\mathbf{X}})^{*}(\theta)}^-{\sim}
&&
(\mathbf{G}\times\mathbf{G}\times\mathbf{X})
  \times_{\mathrm{pr}_{\mathbf{X}}\circ(\mathrm{id}_{\mathbf{G}}\times\lambda)}
  \mathbf{M}
\ar[d]^-{\sim}_-{(\mathrm{id}_{\mathbf{G}}\times\lambda)^{*}(\theta)}
\\
(\mathbf{G}\times\mathbf{G}\times\mathbf{X})
  \times_{\lambda\circ(\mu\times\mathrm{id}_{\mathbf{X}})}
  \mathbf{M}
&&
(\mathbf{G}\times\mathbf{G}\times\mathbf{X})
  \times_{\lambda\circ(\mathrm{id}_{\mathbf{G}}\times\lambda)}
  \mathbf{M}.
}
\]
One therefore sees that giving a structure of \(G\)-equivariant
\(\mathbf{O}_{\mathbf{X}}\)-module on \(\mathbf{M}\) amounts to giving an
isomorphism \(\theta\) of \(\mathbf{O}_{\mathbf{G}\times\mathbf{X}}\)-modules as
above, such that the diagram above is commutative.

(iii) Everything that precedes extends to the case where \(\mathbf{G}\) is
only a \(\widehat{\mathcal{C}}\)-monoid: in this case, giving an action
\(\Lambda\colon\mathbf{G}\times\mathbf{M}\to\mathbf{M}\) lifting \(\lambda\) and
such that each
\(\Lambda_{x}(g)\colon\mathbf{M}_{x}\to\mathbf{M}_{gx}\) is a morphism of
\(\mathbf{O}_{U}\)-modules, amounts to giving a \emph{morphism} \(\theta\) of
\(\mathbf{O}_{\mathbf{G}\times\mathbf{X}}\)-modules as in (ii), such that the
diagram above (without the signs \(\sim\) under the arrows) is commutative,
and such that
\(p_{\mathbf{M}}\circ\theta\circ(\varepsilon_{\mathbf{G}}\times\mathrm{id}_{\mathbf{M}})=\mathrm{id}_{\mathbf{M}}\),
where \(\varepsilon_{\mathbf{G}}\) denotes the unit section of \(\mathbf{G}\)
and \(p_{\mathbf{M}}\) the projection onto \(\mathbf{M}\).
\end{remark}

\par\addvspace{0.55\baselineskip}%
\noindent\textbf{6.3. \(G\)-equivariant morphisms.}\ ---\ \label{I.6.3}%
Let \(\mathbf{Y}\) be a second object of \(\widehat{\mathcal{C}}\), equipped
with an action \(\lambda_{\mathbf{Y}}\colon\mathbf{G}\times\mathbf{Y}\to\mathbf{Y}\)
of \(\mathbf{G}\), and let \(\mathbf{N}\) be a second \(G\)-equivariant
\(\mathbf{O}_{\mathbf{X}}\)-module. One says that a
\(\widehat{\mathcal{C}}\)-morphism \(f\colon\mathbf{Y}\to\mathbf{X}\)
\resp a morphism of \(\mathbf{O}_{\mathbf{X}}\)-modules
\(\phi\colon\mathbf{M}\to\mathbf{N}\) is \emph{\(G\)-equivariant} if it
commutes with the action of \(\mathbf{G}\), \ie if one has, set-theoretically,
\(f(g\cdot y)=g\cdot f(y)\) \resp \(\phi(g\cdot m)=g\cdot\phi(m)\), which
amounts to saying that
\(f\circ\lambda_{\mathbf{Y}}=\lambda_{\mathbf{X}}\circ(\mathrm{id}_{\mathbf{G}}\times f)\)
\resp \(\phi\circ\Lambda_{\mathbf{M}}=\Lambda_{\mathbf{N}}\circ(\mathrm{id}_{\mathbf{G}}\times\phi)\).

One then obtains at once the following lemma:
